\section{Bab 2}
\subsection{Perilaku Demokratis di Era Keterbukaan Informasi}
Nilai-Nilai demokrasi berdasarkan pancasila terdapat pada alinea IV UUD 1945.


\noindent \textbf{Demokrasi: }Demos: Rakyat, Kratos: Kekuasaan $\rightarrow$ Kekuasaan/kekuatan rakyat.


Hakikat Demokrasi $\rightarrow$ 3 pilar $\rightarrow$ Pemerintah (dari, oleh, untuk) rakyat.


\begin{table}[H]
  \centering
  \begin{tabular}{cc}
    \toprule
    \textbf{Demokrasi Langsung} & \textbf{Demokrasi Tidak Langsung} \\ \midrule
    Pemilu (partisipasi langsung) & Memberi kekuasaan pada para \\  \bottomrule
  \end{tabular}
\end{table}


\noindent \textbf{Demokrasi Pancasila: }Pelaksanaan harus atas dasar nilai-nilai pancasila (mengatur antara hubungan masyarakat untuk pemusyawarakatan).


Contoh: Mewujudkan kesejahteraan rakyat berdasarkan keadilan sosial


\noindent \textbf{Demokrasi Liberal} $\rightarrow$ Contoh: Mewujudkan keadilan individual.


\subsubsection{Prinsip-Prinsip Demokrasi}
Menyelesaikan perselisihan, menjamin terselenggara perubahan secara damai, menyelenggarakan pergantian pemimpin.

\subsubsection{Indikator Demokrasi}
\begin{enumerate}
\item Akuntabilitas.
\item Rotasi kekuasaan.
\item Rekrutmen politik terbuka.
\item Pemilu.
\item Hak WNI pasar.
\end{enumerate}

\subsubsection{Keterbukaan Informasi}
Menerima informasi yang akurat, benar dan tidak menyesatkan.


UUD No 14 tahun 2018 $\rightarrow$ Badan publik (eksekusif, legislatif, yuridisktif).


\noindent \textbf{Pasal 28f: }Keterbukaan informasi publik (KIP) UUD No 13 tahun 2008.


\subsubsection{Perilaku Demokratis}
Musyawarah (konsensus), menumbuhkan budaya pasif menghargai kebhinekaan menghormati perbedaan pilihan kesetaraan gender, kebebasan yang bertanggung jawab.


Cara berpikir, bersikap, bertindak berdasarkan persamaan hak dan kewajiban antarai diri dan orang lain.


\noindent \textbf{Kemerdekaan berpendapat: }28 UUD 1945.


\subsubsection{Pentingnya Berkonstitusi}


\textbf{Konstitusi: }Hukum tertinggi yang berisi aturan, norma, prinsip untuk mengatur cara penyelengaraan, pembagian kekuasaan serta hak dan kewajiban WNI.


\noindent \textbf{Jenis Konstitusi Tertulis: }UUD 1945, Tidak tertulis $\rightarrow$ Kebiasaan ketatanegaraan yang tumbuh dipraktikan dalam pemerintahan.


\noindent \textbf{Fungsi Utama: }
\begin{enumerate}
\item Membatasi kekuasaan.
\item Menjamin perlindungan hak asasi manusia.
\item Pedoman dasar.
\end{enumerate}
